\section{Methodology: ThermoMerge}
\label{sec:method}

In this section, we present the mathematical and physical formulation of \textbf{ThermoMerge} (Thermodynamic Model Merging). We begin by establishing the layer-wise parameter merging framework, then formulate the canonical ensemble mapping, rigorously derive the Helmholtz Free Energy Discrepancy objective, and describe the Thermodynamic Annealing Schedule and task-wise thermal coupling.

\subsection{Layer-wise Parameter Merging Formulation}
Let $\theta^{base}$ denote the parameters of a pre-trained base model, and let $\{\theta_k\}_{k=1}^K$ represent a set of $K$ independent task-specific expert networks. These experts are fine-tuned from the common pre-trained ancestor $\theta^{base}$ on distinct task-specific datasets. 

Rather than treating the network parameters as monolithic vectors, we operate at a layer-wise granularity. Let $\theta_l^{base}$ and $\theta_{l,k}$ denote the parameter vectors of layer $l \in \{1, \dots, L\}$ for the base model and expert $k$, respectively. We define the multi-task merged model parameter $\theta_l^{MTL}$ for layer $l$ as:
\begin{equation}
\label{eq:merging_param}
\theta_l^{MTL}(\boldsymbol{\lambda}_l) = \theta_l^{base} + \sum_{k=1}^K \lambda_{l,k} \left( \theta_{l,k} - \theta_l^{base} \right)
\end{equation}
where $\boldsymbol{\lambda}_l = [\lambda_{l,1}, \dots, \lambda_{l,K}]^T \in \mathbb{R}^K$ is a vector of layer-specific, task-specific coupling coefficients. These coefficients govern the degree of parameter-space influence of each expert on the unified model at layer $l$. During unsupervised test-time adaptation, the set of all coupling coefficients $\boldsymbol{\Lambda} = \{\boldsymbol{\lambda}_l\}_{l=1}^L$ serves as the primary optimization parameters.

\subsection{Canonical Ensemble Mapping}
In statistical mechanics, a canonical ensemble describes a physical system in thermal equilibrium with a heat bath at a finite temperature $T$~\cite{gibbs1902elementary}. The probability of the system occupying a microstate $c$ is given by the canonical Boltzmann distribution:
\begin{equation}
P(\text{state } c) = \frac{\exp\left( - \frac{E_c}{k_B T} \right)}{Z}
\end{equation}
where $E_c$ is the energy of state $c$, $k_B$ is the Boltzmann constant (which we set to $1$ without loss of generality), and $Z$ is the partition function that sums over all possible states of the system.

We translate this physics paradigm directly to neural network outputs. Let $f(x; \theta_k) \in \mathbb{R}^{C_k}$ represent the logits output of expert model $k$ on an input $x \in \mathcal{X}_k^{te}$, where $C_k$ is the number of classification classes for task $k$. We map the task classification logits to states in a canonical Boltzmann ensemble, where class-specific logits function as negative microstate energies: $E_c \equiv -f_c(x; \theta_k)$.

Under this physical mapping, the canonical probability distribution for class $c \in \{1, \dots, C_k\}$ under expert $k$ at temperature $T$ is defined as:
\begin{equation}
\label{eq:boltz_expert}
p_c^{(k)}(x; T) = \frac{\exp\left( \frac{f_c(x; \theta_k)}{T} \right)}{Z_k(x; T)}
\end{equation}
where $Z_k(x; T)$ is the expert's canonical partition function, representing the sum over states:
\begin{equation}
\label{eq:partition_expert}
Z_k(x; T) = \sum_{j=1}^{C_k} \exp\left( \frac{f_j(x; \theta_k)}{T} \right)
\end{equation}
The corresponding Helmholtz Free Energy $F_k(x; T)$ of the expert system, which measures the useful work obtainable from a closed thermodynamic system, is defined as:
\begin{equation}
\label{eq:free_energy_expert}
F_k(x; T) = -T \ln Z_k(x; T)
\end{equation}

Similarly, for the merged multi-task network parameterized by $\theta^{MTL}$, the Boltzmann probability distribution for task $k$ at temperature $T$ is written as:
\begin{equation}
\label{eq:boltz_merged}
p_c^{(MTL, k)}(x; T) = \frac{\exp\left( \frac{f_c(x; \theta^{MTL})}{T} \right)}{Z_{MTL, k}(x; T)}
\end{equation}
with the associated multi-task partition function:
\begin{equation}
\label{eq:partition_merged}
Z_{MTL, k}(x; T) = \sum_{j=1}^{C_k} \exp\left( \frac{f_j(x; \theta^{MTL})}{T} \right)
\end{equation}
The Helmholtz Free Energy of the merged model on task $k$ is given by $F_{MTL, k}(x; T) = -T \ln Z_{MTL, k}(x; T)$.

\subsection{Helmholtz Free Energy Discrepancy Minimization}
During test-time adaptation, we have access only to an unlabeled calibration data stream $\mathcal{X}_k^{te}$ for each task $k$. To optimize the coupling coefficients $\boldsymbol{\Lambda}$ without labels, we establish a thermodynamic equilibrium between the merged model and the task experts. We achieve this by minimizing a novel objective: the \textbf{Helmholtz Free Energy Discrepancy (F-Min)} objective, which we define as:
\begin{align}
\label{eq:f_min_obj}
\mathcal{L}(\boldsymbol{\Lambda}, T) = \sum_{k=1}^K \mathbb{E}_{x \in \mathcal{X}_k^{te}} \Big[ T \cdot &\mathcal{D}_{KL} \Big( p^{(k)}(x; T) \parallel \notag \\
&p^{(MTL, k)}(x; T) \Big) \Big]
\end{align}
where $\mathcal{D}_{KL}(\cdot \parallel \cdot)$ is the Kullback-Leibler divergence~\cite{kullback1951information}. Mathematically, this objective is structurally related to temperature-scaled knowledge distillation~\cite{hinton2015distilling}, but here it is derived from first-principles of variational free energy to balance localized expected energy differences and global partition function discrepancies. 

We now present a rigorous derivation of the physical interpretation of this objective. By definition:
\begin{align}
&\mathcal{D}_{KL} \big( p^{(k)} \parallel p^{(MTL, k)} \big) \notag \\
&\quad = \sum_{c=1}^{C_k} p_c^{(k)}(x; T) \ln \left( \frac{p_c^{(k)}(x; T)}{p_c^{(MTL, k)}(x; T)} \right)
\end{align}
Taking the logarithm of the Boltzmann distributions defined in Equations~\eqref{eq:boltz_expert} and~\eqref{eq:boltz_merged}:
\begin{align}
\ln &p_c^{(k)}(x; T) \notag \\
&= \frac{f_c(x; \theta_k)}{T} - \ln Z_k(x; T) \notag \\
&= \frac{f_c(x; \theta_k)}{T} + \frac{F_k(x; T)}{T} \\[0.5em]
\ln &p_c^{(MTL, k)}(x; T) \notag \\
&= \frac{f_c(x; \theta^{MTL})}{T} - \ln Z_{MTL, k}(x; T) \notag \\
&= \frac{f_c(x; \theta^{MTL})}{T} + \frac{F_{MTL, k}(x; T)}{T}
\end{align}
Subtracting these two expressions:
\begin{align}
\ln \left( \frac{p_c^{(k)}(x; T)}{p_c^{(MTL, k)}(x; T)} \right) &= \frac{f_c(x; \theta_k) - f_c(x; \theta^{MTL})}{T} \notag \\
&\quad + \frac{F_k(x; T) - F_{MTL, k}(x; T)}{T}
\end{align}
Multiplying both sides by the temperature $T$:
\begin{align}
T \ln \left( \frac{p_c^{(k)}(x; T)}{p_c^{(MTL, k)}(x; T)} \right) &= f_c(x; \theta_k) - f_c(x; \theta^{MTL}) \notag \\
&\quad + F_k(x; T) - F_{MTL, k}(x; T)
\end{align}
Finally, taking the expectation over all classes $c$ under the expert distribution $p^{(k)}(x; T)$:
\begin{align}
\label{eq:derivation}
&T \cdot \mathcal{D}_{KL} \big( p^{(k)} \parallel p^{(MTL, k)} \big) \notag \\
&\quad = \sum_{c} p_c^{(k)} ( f_c(x; \theta_k) - f_c(x; \theta^{MTL}) ) \notag \\
&\quad + ( F_k(x; T) - F_{MTL, k}(x; T) )
\end{align}

Equation~\eqref{eq:derivation} reveals a profound physical duality. Here, we clarify a crucial distinction in thermodynamic terminology. The scaled KL divergence $T \cdot \mathcal{D}_{KL}(p^{(k)} \parallel p^{(MTL, k)})$ represents the difference between the \emph{variational free energy} of the expert distribution $p^{(k)}$ evaluated under the merged model's energy states and the true \emph{equilibrium free energy} of the merged model $F_{MTL, k}(x; T)$, rather than a simple difference between their equilibrium free energies ($F_k - F_{MTL, k}$). It is this variational energy gap that drives parameter adaptation.

Specifically, the right-hand side of Equation~\eqref{eq:derivation} comprises two terms:
\begin{enumerate}
    \item \textbf{Expected Energy Difference:} $\sum_{c=1}^{C_k} p_c^{(k)}(x; T) \left( f_c(x; \theta_k) - f_c(x; \theta^{MTL}) \right)$ measures the average energy difference between the expert's microstates and those of the merged multi-task model, weighted by the expert's thermal probability distribution.
    \item \textbf{Helmholtz Free Energy Difference:} $F_k(x; T) - F_{MTL, k}(x; T)$ captures the global thermodynamic state discrepancy between the expert and merged model partition functions.
\end{enumerate}

Thus, the Helmholtz Free Energy Discrepancy (F-Min) objective naturally balances localized energy minimization with global thermodynamic state matching. This physical balancing act acts as a powerful regularizer, anchoring the optimized coefficients $\boldsymbol{\Lambda}$ and protecting the merged model from the transductive overfitting that plagues standard entropy-based test-time adaptation. We provide a complete, step-by-step algebraic expansion and proof of this first-principles derivation in Appendix~\ref{sec:math_derivation}.

\subsection{Thermodynamic Annealing Schedule (TAS)}
Test-time adaptation landscapes are notoriously non-convex and littered with poor local minima, which are direct artifacts of linear parameter merging across disjoint loss basins. To bypass these optimization barriers, we introduce a simulated \textbf{Thermodynamic Annealing Schedule (TAS)}~\cite{kirkpatrick1983optimization}, drawing inspiration from classical physical cooling processes and annealed statistical sampling methods~\cite{neal2001annealed}. We parameterize the system's global temperature $T(t)$ as a dynamic, decaying function of the optimization step $t \in \{1, \dots, N_{steps}\}$:
\begin{equation}
\label{eq:annealing}
T(t) = T_{end} + \left( T_{start} - T_{end} \right) \cdot \exp\left( -\beta \cdot t \right)
\end{equation}
where $T_{start}$ is the hot initial temperature, $T_{end}$ is the target inference temperature, and $\beta > 0$ is the thermal cooling rate.

During the initial phase of optimization (high $T(t)$), the Boltzmann distributions are highly diffuse, which physically corresponds to high thermal excitation. In this state, the free energy landscapes are mathematically flattened, bridging disjoint basins of attraction and allowing the trainable coupling parameters $\boldsymbol{\Lambda}$ to slide smoothly over high-loss Euclidean ridges without getting trapped in frustrated local minima. As the system cools down (decreasing $T(t)$), the probability distributions crystallize around specialized classification boundaries, locking in high-performance multi-task representations. In our experiments, we set $T_{start} = 5.0$, $T_{end} = 1.0$, and $\beta = 0.05$. We provide a deeper physical conceptual grounding of this simulated thermalization process and its connection to spin glass theory, multi-task parameter frustration, and replica symmetry breaking in Appendix~\ref{sec:spin_glass_details}.

\subsection{Task-wise Thermal Coupling}
\label{sec:task_coupling}
To further align our physical framework with multi-task learning scenarios where distinct downstream tasks exhibit vastly different levels of baseline confidence and difficulty, we introduce trainable, task-wise local temperatures $T_k(t)$. Physically, each downstream task $k \in \{1,\dots,K\}$ is modeled as a specialized thermodynamic subsystem in contact with its own local thermal bath, characterized by a trainable local thermal capacity $\tau_k > 0$.

Specifically, we couple the global physical temperature $T(t)$ governed by the Thermodynamic Annealing Schedule (Equation~\eqref{eq:annealing}) with task-specific thermal scaling factors to define the local temperature $T_k(t)$ of task $k$:
\begin{equation}
\label{eq:thermo_task}
T_k(t) = \tau_k \cdot T(t)
\end{equation}
where $\tau_k > 0$ is a trainable local task coupling coefficient initialized to $1.0$ (corresponding to a baseline heat capacity). 

To prevent division by zero or extreme scaling during the gradient descent optimization, we apply a strict physical stabilization constraint by clamping the thermal capacity parameters to a bounded interval: $\tau_k \in [0.2, 5.0]$. This specific clamping range was chosen based on both empirical stability and numerical constraints. Under standard floating-point precision, local temperatures below $0.2$ lead to extreme logit amplification, causing exponential values to overflow and trigger NaN gradients during backpropagation. Conversely, local temperatures above $5.0$ excessively dilute the output logits, flattening the prediction probabilities into uniform noise and destroying the gradient signal required for adaptation. This bounded interval thus guarantees numerical safety while preserving sufficient capacity for local thermal adjustments. 

Under this Task-wise Thermal Coupling, the Boltzmann distributions for task $k$ are scaled by its own localized, dynamic temperature $T_k(t)$:
\begin{align}
p_c^{(k)}(x; T_k(t)) &= \frac{\exp\left( \frac{f_c(x; \theta_k)}{T_k(t)} \right)}{Z_k(x; T_k(t))} \\
p_c^{(MTL, k)}(x; T_k(t)) &= \frac{\exp\left( \frac{f_c(x; \theta^{MTL})}{T_k(t)} \right)}{Z_{MTL, k}(x; T_k(t))}
\end{align}
The associated Free Energy Discrepancy is minimized over both layer-wise merging coefficients $\boldsymbol{\Lambda}$ and task-wise thermal coupling capacities $\boldsymbol{\tau} = [\tau_1,\dots,\tau_K]^T$:
\begin{equation}
\{\boldsymbol{\Lambda}^*, \boldsymbol{\tau}^*\} = \arg\min_{\boldsymbol{\Lambda}, \boldsymbol{\tau}} \mathcal{L}(\boldsymbol{\Lambda}, \boldsymbol{\tau})
\end{equation}

By scaling the Boltzmann distributions and Helmholtz Free Energy of each task $k$ using its own local temperature $T_k(t)$, we ensure that simple tasks (which quickly converge and require sharp, low-temperature boundaries) and complex tasks (which require soft, high-temperature boundaries to prevent collapse) can discover their own stable localized thermal equilibrium. During final evaluation, $T_k(1.0) = \tau_k \cdot 1.0 = \tau_k$ represents the crystallized optimal thermal scaling parameter for output predictions on task $k$.

We clarify a subtle but mathematically important property of this temperature scaling at test-time inference. After test-time adaptation is complete, the final classification decision for a target input $x$ is obtained via an $\arg\max$ operator over the task's output logits: $\hat{c} = \arg\max_{j} \left( f_j(x; \theta^{MTL}) / T_k \right)$. Since the crystallized temperature $T_k = \tau_k$ is strictly positive ($\tau_k \in [0.2, 5.0]$), dividing the logits by a constant positive scalar does not alter their rank-ordering. Consequently, the final classification accuracy on any individual task is mathematically invariant to the value of $\tau_k$ during evaluation. Instead, the task-wise capacities $\tau_k$ play a crucial, dynamic regularization role \emph{strictly during the optimization/adaptation phase} by shaping the gradients of the Free Energy loss and controlling the rate of parameter convergence.

Importantly, this Task-wise Thermal Coupling operates entirely within the output-space logit Boltzmann ensemble, strictly preserving the standard parameter-space model merging formulation of Equation~\eqref{eq:merging_param} without introducing any ad-hoc weight scaling heuristics in the internal network layers.
